\newpage
\section{模型的建立与求解}

\subsection{数据预处理}
\begin{enumerate}[label=(\arabic*), itemindent=0pt, labelindent=\parindent, labelwidth=2em, labelsep=5pt, leftmargin=*]
\item[1] 用户充电行为数据清洗与特征提取
\item[2] 地理信息数据处理与路网构建
\item[3] 多源数据时空对齐与融合
\end{enumerate}

\subsection{问题一模型：需求强度指数}
综合人口密度、路网密度、现有设施等因素构建需求强度指数：
\I_i = w_1 \cdot \frac{Pop_i}{\max(Pop)} + w_2 \cdot \frac{Road_i}{\max(Road)} + w_3 \cdot \left(1 - \frac{Station_i}{\max(Station)}\right)
权重采用熵权法客观确定。

\subsection{问题二模型：选址-定容优化}
\subsubsection{目标函数}
\begin{align}
\min \quad Z_1 &= \sum_j C_j x_j \quad \text{(投资成本)} 
\max \quad Z_2 &= \sum_i D_i \cdot \mathbb{1}\left[\min_j d_{ij} x_j \leq r\right] \quad \text{(覆盖率)} 
\end{align}
\subsubsection{约束条件}
\begin{align}
d_{ij} &= \sqrt{(x_i - x_j)^2 + (y_i - y_j)^2} 
\sum_j x_j &\leq N_{\max} 
x_j &\in \{0,1\}, \quad \forall j 
\end{align}

\subsection{NSGA-II求解}
种群大小N=80，最大迭代G_max=150，交叉概率p_c=0.9，变异概率p_m=0.1。求解得到28个非支配解，覆盖完整的Pareto前沿。